\documentclass[11pt,a4paper]{article}

\usepackage[margin=18mm]{geometry}
\usepackage{amsmath,amssymb}
\usepackage{array}
\usepackage{booktabs}
\usepackage[T1]{fontenc}
\usepackage{lmodern}
\usepackage{microtype}
\usepackage[hidelinks]{hyperref}

\setlength{\parindent}{0pt}
\setlength{\parskip}{0.65em}

\newcommand{\ket}[1]{\left|#1\right\rangle}
\newcommand{\bra}[1]{\left\langle#1\right|}
\newcommand{\braket}[2]{\left\langle #1 \middle| #2 \right\rangle}

\title{\textbf{Quantum Computing Linear Algebra - Quick Reference}}
\author{}
\date{}

\begin{document}

\maketitle
\vspace{-1.5em}

This reference keeps the ordinary linear-algebra representation next to the
quantum notation. The main question to ask is: \emph{what objects are being
combined, and what kind of object should the operation produce?}

\section{Column vectors, row vectors, transpose, and adjoint}

A state vector is normally represented as a \textbf{column vector}:

\[
a =
\begin{bmatrix}
a_0\\
a_1
\end{bmatrix},
\qquad
\text{shape } 2\times1.
\]

Its ordinary transpose is a \textbf{row vector}:

\[
a^T =
\begin{bmatrix}
a_0 & a_1
\end{bmatrix},
\qquad
\text{shape } 1\times2.
\]

Quantum states may contain complex numbers. Therefore quantum notation normally
uses the \textbf{conjugate transpose} (also called the Hermitian adjoint):

\[
a^\dagger =
\begin{bmatrix}
a_0^* & a_1^*
\end{bmatrix}.
\]

Here $^*$ means complex conjugation. If all entries are real, then
$a^\dagger=a^T$.

In Dirac notation:

\[
\ket{a}
\quad\longleftrightarrow\quad
\begin{bmatrix}
a_0\\a_1
\end{bmatrix},
\qquad
\bra{a}
\quad\longleftrightarrow\quad
\begin{bmatrix}
a_0^* & a_1^*
\end{bmatrix}.
\]

\section{Inner product}

\textbf{Input:} row vector $\times$ column vector

\textbf{Output:} scalar

\textbf{Quantum notation:}
\[
\braket{a}{b}.
\]

For example,

\[
\ket{a} =
\begin{bmatrix}
1\\0
\end{bmatrix},
\qquad
\ket{b} =
\frac{1}{\sqrt{2}}
\begin{bmatrix}
1\\1
\end{bmatrix}.
\]

Then

\[
\braket{a}{b}
=
\begin{bmatrix}
1 & 0
\end{bmatrix}
\frac{1}{\sqrt{2}}
\begin{bmatrix}
1\\1
\end{bmatrix}
=
\frac{1}{\sqrt{2}}
(1\cdot1+0\cdot1)
=
\frac{1}{\sqrt{2}}.
\]

The corresponding squared magnitude is

\[
\left|\braket{a}{b}\right|^2
=
\left|\frac{1}{\sqrt{2}}\right|^2
=
\frac{1}{2}.
\]

\textbf{Quantum use:} state overlap and measurement probabilities.

\section{Outer product}

\textbf{Input:} column vector $\times$ row vector

\textbf{Output:} matrix

For

\[
a =
\begin{bmatrix}
a_0\\a_1
\end{bmatrix},
\qquad
b =
\begin{bmatrix}
b_0\\b_1
\end{bmatrix},
\]

the real-valued outer product is

\[
ab^T
=
\begin{bmatrix}
a_0\\a_1
\end{bmatrix}
\begin{bmatrix}
b_0 & b_1
\end{bmatrix}
=
\begin{bmatrix}
a_0b_0 & a_0b_1\\
a_1b_0 & a_1b_1
\end{bmatrix}.
\]

In quantum notation we normally use the bra, so

\[
\ket{a}\bra{b}
\]

uses the conjugate transpose of $b$.

A simple example is

\[
\ket{0}\bra{0}
=
\begin{bmatrix}
1\\0
\end{bmatrix}
\begin{bmatrix}
1 & 0
\end{bmatrix}
=
\begin{bmatrix}
1 & 0\\
0 & 0
\end{bmatrix}.
\]

\textbf{Quantum use:} projectors and density matrices. For a pure state,

\[
\rho = \ket{\psi}\bra{\psi}.
\]

\section{Kronecker / tensor product}

\textbf{Input:} state vector $\otimes$ state vector

\textbf{Output:} larger state vector

For

\[
a =
\begin{bmatrix}
a_0\\a_1
\end{bmatrix},
\qquad
b =
\begin{bmatrix}
b_0\\b_1
\end{bmatrix},
\]

the Kronecker product is

\[
a\otimes b
=
\begin{bmatrix}
a_0 b\\
a_1 b
\end{bmatrix}
=
\begin{bmatrix}
a_0b_0\\
a_0b_1\\
a_1b_0\\
a_1b_1
\end{bmatrix}.
\]

For example,

\[
\ket{0}\otimes\ket{1}
=
\begin{bmatrix}
1\\0
\end{bmatrix}
\otimes
\begin{bmatrix}
0\\1
\end{bmatrix}
=
\begin{bmatrix}
1\cdot0\\
1\cdot1\\
0\cdot0\\
0\cdot1
\end{bmatrix}
=
\begin{bmatrix}
0\\1\\0\\0
\end{bmatrix}
=
\ket{01}.
\]

For two qubits, the computational basis is

\[
\ket{00}=
\begin{bmatrix}1\\0\\0\\0\end{bmatrix},
\quad
\ket{01}=
\begin{bmatrix}0\\1\\0\\0\end{bmatrix},
\quad
\ket{10}=
\begin{bmatrix}0\\0\\1\\0\end{bmatrix},
\quad
\ket{11}=
\begin{bmatrix}0\\0\\0\\1\end{bmatrix}.
\]

Dimension growth is

\[
2\otimes2=4,
\qquad
2\otimes2\otimes2=8,
\qquad
n\text{ qubits}\longrightarrow2^n\text{ amplitudes}.
\]

\textbf{Key distinction:} the outer product produces a matrix; the Kronecker
product composes systems into a larger state vector.

\section{Matrix-vector product}

\textbf{Input:} operator matrix $\times$ state vector

\textbf{Output:} new state vector

\textbf{Quantum notation:}
\[
U\ket{\psi}.
\]

Consider the Hadamard gate:

\[
H=
\frac{1}{\sqrt{2}}
\begin{bmatrix}
1 & 1\\
1 & -1
\end{bmatrix},
\qquad
\ket{0}=
\begin{bmatrix}
1\\0
\end{bmatrix}.
\]

Apply the matrix to the vector:

\[
H\ket{0}
=
\frac{1}{\sqrt{2}}
\begin{bmatrix}
1 & 1\\
1 & -1
\end{bmatrix}
\begin{bmatrix}
1\\0
\end{bmatrix}.
\]

Calculate each output row:

\[
=
\frac{1}{\sqrt{2}}
\begin{bmatrix}
1\cdot1+1\cdot0\\
1\cdot1+(-1)\cdot0
\end{bmatrix}
=
\frac{1}{\sqrt{2}}
\begin{bmatrix}
1\\1
\end{bmatrix}.
\]

Using the basis vectors,

\[
\frac{1}{\sqrt{2}}
\begin{bmatrix}
1\\1
\end{bmatrix}
=
\frac{1}{\sqrt{2}}
\left(
\begin{bmatrix}
1\\0
\end{bmatrix}
+
\begin{bmatrix}
0\\1
\end{bmatrix}
\right)
=
\frac{\ket{0}+\ket{1}}{\sqrt{2}}.
\]

So the familiar circuit operation

\[
\ket{0}\xrightarrow{H}
\frac{\ket{0}+\ket{1}}{\sqrt{2}}
\]

is ordinary matrix-vector multiplication.

\section{Matrix-matrix product}

\textbf{Input:} operator matrix $\times$ operator matrix

\textbf{Output:} composed operator matrix

If $U_1$ acts first and $U_2$ acts second,

\[
\ket{\psi_1}=U_1\ket{\psi},
\]

\[
\ket{\psi_2}=U_2\ket{\psi_1}
            =U_2U_1\ket{\psi}.
\]

The \textbf{rightmost operator acts first}.

Our first experiment used the fact that

\[
H^2=HH=I.
\]

Start with

\[
HH
=
\left(
\frac{1}{\sqrt{2}}
\begin{bmatrix}
1 & 1\\
1 & -1
\end{bmatrix}
\right)
\left(
\frac{1}{\sqrt{2}}
\begin{bmatrix}
1 & 1\\
1 & -1
\end{bmatrix}
\right).
\]

Combine the scalar factors:

\[
=
\frac{1}{2}
\begin{bmatrix}
1 & 1\\
1 & -1
\end{bmatrix}
\begin{bmatrix}
1 & 1\\
1 & -1
\end{bmatrix}.
\]

Calculate each matrix entry as row $\cdot$ column:

\[
=
\frac{1}{2}
\begin{bmatrix}
1\cdot1+1\cdot1 &
1\cdot1+1\cdot(-1)
\\[4pt]
1\cdot1+(-1)\cdot1 &
1\cdot1+(-1)\cdot(-1)
\end{bmatrix}.
\]

Therefore,

\[
=
\frac{1}{2}
\begin{bmatrix}
2 & 0\\
0 & 2
\end{bmatrix}
=
\begin{bmatrix}
1 & 0\\
0 & 1
\end{bmatrix}
=
I.
\]

Applying this composed operator to $\ket{0}$ gives

\[
HH\ket{0}
=
I\ket{0}
=
\ket{0},
\]

which is exactly what we observed experimentally with 1000 measurements of
\texttt{0}.

\section{Bell-state connection}

The Bell state from the current lab is

\[
\ket{\Phi^+}
=
\frac{\ket{00}+\ket{11}}{\sqrt{2}}.
\]

Using the computational basis,

\[
=
\frac{1}{\sqrt{2}}
\left(
\begin{bmatrix}
1\\0\\0\\0
\end{bmatrix}
+
\begin{bmatrix}
0\\0\\0\\1
\end{bmatrix}
\right)
=
\frac{1}{\sqrt{2}}
\begin{bmatrix}
1\\0\\0\\1
\end{bmatrix}.
\]

After applying $Z$ to the appropriate qubit, our simulator showed

\[
\ket{\Phi^-}
=
\frac{\ket{00}-\ket{11}}{\sqrt{2}}
=
\frac{1}{\sqrt{2}}
\begin{bmatrix}
1\\0\\0\\-1
\end{bmatrix}.
\]

The signs differ, but direct computational-basis measurement still gives

\[
P(00)=\frac{1}{2},
\qquad
P(11)=\frac{1}{2}.
\]

The relative phase becomes observable after further gate operations create
interference.

\section{At-a-glance summary}

\begin{center}
\begin{tabular}{@{}llll@{}}
\toprule
Operation & Typical notation & Result & Quantum role\\
\midrule
Inner product & $\braket{a}{b}$ & scalar & overlap / probability\\
Outer product & $\ket{a}\bra{b}$ & matrix & projector / density matrix\\
Kronecker product & $\ket{a}\otimes\ket{b}$ & larger vector & combine systems\\
Matrix-vector & $U\ket{\psi}$ & vector & apply gate\\
Matrix-matrix & $U_2U_1$ & matrix & compose gates\\
\bottomrule
\end{tabular}
\end{center}

\end{document}
